% Orthogonalität – bank/orthogonalitaet/ (zusammenbau v0.4, Heft)
\einheitenkopf[t9]{Orthogonalität}
\setcounter{aufgabe}{80}

% Anlauf (ohne Prüfkennung)
\begin{aufgabe}{Rechte Winkel nachweisen – Anlauf}
\begin{teile}
\teil Das Dreieck $ABC$ soll bei $B$ einen rechten Winkel haben. Welche Vektoren prüfst du? \\ \kreuz{$\overrightarrow{BA}$ und $\overrightarrow{BC}$} \\ \kreuz{$\overrightarrow{AB}$ und $\overrightarrow{AC}$} \\ \kreuz{$\overrightarrow{CA}$ und $\overrightarrow{CB}$}
\teil Gegeben sind $A(2 | 1 | 3)$, $B(5 | 3 | 3)$, $C(0 | 4 | 3)$. Weise nach, dass das Dreieck bei $A$ einen rechten Winkel hat.
\teil Gegeben sind $P(4 | 0 | 0)$, $Q(0 | 3 | 5)$ und der Ursprung $O$. Begründe ohne Skalarprodukt, dass das Dreieck $OPQ$ bei $O$ rechtwinklig ist.
\end{teile}
\end{aufgabe}

\begin{aufgabe}{Rechte Winkel nachweisen – Prüfungsaufgaben}
\begin{teile}
\teil Das Dreieck $OAB$ mit $O(0 | 0 | 0)$, $A(5 | 12 | 0)$ und $B(0 | 0 | 13)$ ist gleichschenklig. Weise nach, dass es bei $O$ rechtwinklig ist und dass beide Katheten die Länge 13 haben \hfill (Abitur 2026 GK).
\teil Ein Holzkörper hat die Eckpunkte $A(0 | 0 | 0)$, $B(9 | 0 | 0)$, $C(9 | 6 | 0)$, $D(0 | 6 | 0)$ und $E(0 | 6 | 4)$; er ist eine Pyramide über dem Rechteck $ABCD$ mit der Spitze $E$ senkrecht über $D$ (1 LE = 1 cm). Weise nach, dass das Dreieck $BCE$ rechtwinklig ist, und berechne den Inhalt der gesamten Oberfläche \hfill (Abitur 2021 GK).
\teil Ein Holzkörper hat die Eckpunkte $A(0 | 0 | 0)$, $B(7 | 0 | 0)$, $C(7 | 7 | 0)$, $D(0 | 7 | 0)$ und $E(7 | 0 | 5)$; er ist eine Pyramide über dem Quadrat $ABCD$ mit der Spitze $E$ senkrecht über $B$ (1 LE = 1 cm). Weise nach, dass das Dreieck $CDE$ bei $C$ rechtwinklig ist, gib die Kathetenlängen an und berechne den Inhalt der Oberfläche des Körpers \hfill (Abitur 2021 GK).
\teil Gegeben sind $A(0 | 0 | 0)$, $B(4 | 3 | 5)$ und $C_t(3t | -4t | 0)$ mit $t > 0$. Weise nach, dass das Dreieck $ABC_t$ für jedes $t$ bei $A$ einen rechten Winkel hat \hfill (Abitur 2026 GK).
\teil Gegeben sind $A(1 | 6 | 2)$, $B(4 | 6 | 0)$ und $C_t(4 + 2t | 6 | 3t)$ mit $t \ne 0$. Weise nach, dass jedes Dreieck $ABC_t$ bei $B$ einen rechten Winkel hat \hfill (Abitur 2018 GK).
\teil Ein gerades Prisma $ABCDEF$ hat die Grundfläche $ABC$ mit $A(0 | -6 | 0)$, $B(\sqrt{7} | 0 | 0)$ und $C(0 | 6 | 0)$. Weise nach, dass das Dreieck $ABC$ bei $B$ nicht rechtwinklig ist \hfill (Abitur 2019 GK).
\teil Ein gerades Prisma hat die Grundfläche $ABC$ mit $A(0 | -2 | 0)$, $B(\sqrt{5} | 0 | 0)$ und $C(0 | 3 | 0)$. Weise nach, dass das Dreieck $ABC$ bei $B$ nicht rechtwinklig ist \hfill (Abitur 2019 GK).
\end{teile}
\end{aufgabe}

% Anlauf (ohne Prüfkennung)
\begin{aufgabe}{Rechte Winkel rückwärts – Anlauf}
\begin{teile}
\teil Aufgabe: „Bestimme $a$ so, dass das Dreieck $OAB$ bei $B$ rechtwinklig ist.“ Was ist zu tun? \\ \kreuz{null zeigen – für jeden Wert des Parameters} \\ \kreuz{null setzen – ein Wert ist gesucht, Gleichung lösen} \\ \kreuz{null zeigen – ohne Parameter, einsetzen und ausrechnen}
\teil Gegeben sind $A(1 | 2 | 3)$, $B(4 | 3 | 3)$ und $C(2 | t | 5)$. Bestimme $t$ so, dass das Dreieck $ABC$ bei $A$ rechtwinklig ist.
\teil Ein Quader hat die Eckpunkte $O(0 | 0 | 0)$, $A(12 | 0 | 0)$, $B(12 | 9 | 0)$, $C(0 | 9 | 0)$ und $E(12 | 0 | h)$ mit $h > 0$; $G$ liegt über $C$. Die Raumdiagonalen $AG$ und $CE$ stehen senkrecht zueinander. Bestimme $h$ sowie Volumen und Oberflächeninhalt des Quaders.
\end{teile}
\end{aufgabe}

\begin{aufgabe}{Rechte Winkel rückwärts – Prüfungsaufgaben}
\begin{teile}
\teil Gegeben sind der Ursprung $O$, $A(4 | 0 | a)$ und $B(2 | 2 | 3)$. Bestimme $a$ so, dass das Dreieck $OAB$ bei $B$ rechtwinklig ist \hfill (Abitur 2021 GK).
\teil Der Ursprung $O$, $P(4 | 0 | 12)$ und $Q_t(2 | t | 8)$ bilden ein Dreieck. Bestimme alle Werte von $t$, für die das Dreieck bei $Q_t$ einen rechten Winkel hat \hfill (Abitur 2024 GK).
\teil Gegeben sind $M(5 | 1 | 2)$, $F(0 | 0 | 2)$ und $S_t(t | 0 | 0)$. Bestimme alle $t$, für die das Dreieck $MFS_t$ bei $S_t$ einen rechten Winkel hat \hfill (Abitur 2024 GK).
\teil Gegeben sind $A(1 | 2 | 0)$ und $B(4 | 4 | 2)$. Für den Punkt $C$ gilt: Seine $z$-Koordinate ist $5$, die Geraden $AB$ und $AC$ verlaufen senkrecht zueinander, und $C$ liegt in der Ebene $2x + y = 0$. Bestimme die Koordinaten von $C$ \hfill (Abitur 2026 GK).
\teil Gegeben sind $A(0 | 1 | 2)$ und $B(1 | 4 | 1)$. Für den Punkt $C$ gilt: Seine $z$-Koordinate ist $-1$, die Geraden $AB$ und $AC$ verlaufen senkrecht zueinander, und $C$ liegt in der Ebene $x + y = 2$. Bestimme die Koordinaten von $C$ \hfill (Abitur 2026 GK).
\end{teile}
\end{aufgabe}

% Anlauf (ohne Prüfkennung)
\begin{aufgabe}{Senkrecht zu Geraden und Ebenen – Anlauf}
\begin{teile}
\teil Aufgabe: „Zeige, dass die Gerade $h$ senkrecht auf der Ebene $E$ steht.“ Welche Prüfregel? \\ \kreuz{zwei Richtungen – Skalarprodukt null} \\ \kreuz{Gerade gegen Ebene – Richtungsvektor kollinear zum Normalenvektor} \\ \kreuz{Ebene gegen Ebene – Skalarprodukt der Normalenvektoren null}
\teil Gegeben sind $E\colon x + y + z = 1$ und $F\colon 2x - ty + 3z = 0$. Bestimme $t$ so, dass $E$ und $F$ senkrecht zueinander stehen.
\end{teile}
\end{aufgabe}

\begin{aufgabe}{Senkrecht zu Geraden und Ebenen – Prüfungsaufgaben}
\begin{teile}
\teil Ein Prisma $ABCDEF$ hat die Grundfläche $ABC$ in der Ebene $L\colon 2x - y + 2z = 6$ mit $A(3 | 0 | 0)$, $B(1 | 2 | 3)$, $C(0 | -2 | 2)$; die Seitenkante $BE$ führt zu $E(5 | 0 | 7)$. Begründe, dass das Prisma gerade ist \hfill (Abitur 2020 GK).
\teil Gegeben sind $E\colon \vec x = (2 | 0 | 1) + r \cdot (2 | 1 | 2) + s \cdot (1 | 0 | -1)$ und der Vektor $\vec n = (1 | -4 | 1)$. Weise nach, dass $\vec n$ senkrecht zu $E$ steht \hfill (Abitur 2025 GK).
\teil Ein Kirchturmdach hat unter anderem die Eckpunkte $B(9 | 0 | 0)$, $C(9 | 9 | 0)$, $G(6 | 9 | 5)$ und $F(9 | 6 | 5)$ (1 LE = 1 m). Die Ebene $L$ enthält $C$, $G$ und $F$; die Ebene $N$ ist parallel zu $L$ und enthält $B$. Weise nach, dass $\vec n = (5 | 5 | 3)$ ein Normalenvektor von $N$ ist, und gib eine Koordinatengleichung von $N$ an \hfill (Abitur 2022 GK).
\teil Gegeben sind $E\colon \vec x = (0 | 1 | 3) + r \cdot (0 | 2 | 1) + s \cdot (1 | 1 | 0)$ und $\vec n = (1 | -1 | 2)$. Weise nach, dass $\vec n$ senkrecht zu $E$ steht \hfill (Abitur 2025 GK).
\teil Gegeben ist $g\colon \vec x = (5 | 2 | -1) + s \cdot (-3 | 0 | 2)$; die Gerade $h$ verläuft parallel zur $y$-Achse und schneidet $g$ im Punkt $(5 | 2 | -1)$. Untersuche, ob $g$ und $h$ senkrecht zueinander verlaufen \hfill (Abitur 2025 GK).
\teil Gegeben ist $g\colon \vec x = (1 | 4 | 2) + t \cdot (2 | -3 | 0)$; die Gerade $h$ verläuft parallel zur $x$-Achse durch $(1 | 4 | 2)$. Untersuche, ob $g$ und $h$ senkrecht zueinander verlaufen \hfill (Abitur 2025 GK).
\teil Gegeben sind $E\colon 4x - 3y + z = 2$ und $F\colon x + 2y + kz = 5$. Bestimme $k$ so, dass $F$ senkrecht zu $E$ steht \hfill (Abitur 2022 GK).
\teil Gegeben sind $E\colon 2y - 5z = 0$ und $F\colon 3x + ay + 2z = 7$. Bestimme $a$ so, dass $F$ senkrecht zu $E$ steht \hfill (Abitur 2022 GK).
\teil Gegeben sind $E\colon x + 2y - z = 3$ und $F\colon 3x - y + 2z = 1$; ihre Schnittgerade hat den Richtungsvektor $(3 | -5 | -7)$. Gib eine Gleichung einer Ebene $H$ an, die senkrecht zu $E$ und zu $F$ steht \hfill (Abitur 2021 GK).
\teil Gegeben sind $E\colon 2x - y + z = 4$ und $F\colon x + y - 3z = 2$; ihre Schnittgerade hat den Richtungsvektor $(2 | 7 | 3)$. Gib eine Gleichung einer Ebene $H$ an, die senkrecht zu $E$ und zu $F$ steht \hfill (Abitur 2021 GK).
\end{teile}
\end{aufgabe}

% Anlauf (ohne Prüfkennung)
\begin{aufgabe}{Flächengleichung des flächenkleinsten Dreiecks begründen – Anlauf}
\begin{teile}
\teil Gegeben sind $R(2 | 0 | 1)$, $S(2 | 0 | 5)$ und die Gerade $t\colon \vec x = (0 | 1 | 3) + \lambda \cdot (1 | 2 | 0)$. Jeder Punkt $T_\lambda$ von $t$ hat von $R$ und $S$ denselben Abstand. Begründe, dass das Dreieck $RST_\lambda$ mit dem kleinsten Flächeninhalt zu dem $\lambda$ gehört, das $(\lambda - 2 | 1 + 2\lambda | 0) \circ (1 | 2 | 0) = 0$ löst, und gib dieses $\lambda$ an.
\end{teile}
\end{aufgabe}

\begin{aufgabe}{Flächengleichung des flächenkleinsten Dreiecks begründen – Prüfungsaufgaben}
\begin{teile}
\teil Gegeben sind $R(0 | 1 | 3)$, $S(0 | 5 | 3)$ und die Gerade $t\colon \vec x = (1 | 3 | 0) + \lambda \cdot (2 | 0 | 1)$. Jeder Punkt $T_\lambda$ von $t$ hat von $R$ und $S$ denselben Abstand. Unter den Dreiecken $RST_\lambda$ hat eines den kleinsten Flächeninhalt. Begründe, dass der zugehörige Wert von $\lambda$ die Gleichung $(1 + 2\lambda | 0 | \lambda - 3) \circ (2 | 0 | 1) = 0$ löst \hfill (Abitur 2020 GK).
\end{teile}
\end{aufgabe}
